{
\footnotesize
\setlength{\tabcolsep}{3pt}
\begin{longtable}{|>{\centering\arraybackslash}m{0.75cm}|m{2cm}|m{8cm}|>{\centering\arraybackslash}m{1.75cm}|}
\caption{Kebutuhan Fungsional Sistem}
\label{tab:kebutuhan_fungsional} \\
\hline
\multicolumn{1}{|>{\centering\arraybackslash}m{0.75cm}|}{\textbf{ID}} & \multicolumn{1}{>{\centering\arraybackslash}m{2cm}|}{\textbf{Kebutuhan Fungsional}} & \multicolumn{1}{>{\centering\arraybackslash}m{8cm}|}{\textbf{Deskripsi}} & \multicolumn{1}{>{\centering\arraybackslash}m{1.75cm}|}{\textbf{Tujuan}} \\
\hline
\endfirsthead
\hline
\multicolumn{1}{|>{\centering\arraybackslash}m{0.75cm}|}{\textbf{ID}} & \multicolumn{1}{>{\centering\arraybackslash}m{2cm}|}{\textbf{Kebutuhan Fungsional}} & \multicolumn{1}{>{\centering\arraybackslash}m{8cm}|}{\textbf{Deskripsi}} & \multicolumn{1}{>{\centering\arraybackslash}m{1.75cm}|}{\textbf{Tujuan}} \\
\hline
\endhead
\endfoot
\hline
\endlastfoot

KF-01 & Manajemen Fitur Terpusat & Sistem menyediakan \textit{feature registry} terpusat berbasis Feast dengan \textit{file-based registry} pada MinIO untuk mendefinisikan, mengelola, dan melakukan \textit{versioning} fitur. \textit{Registry} menyimpan metadata seperti deskripsi, tipe data, \textit{owner}, dan dependensi, serta dikelola di dalam repositori Gitea agar mendukung \textit{workflow} GitOps dan penelusuran perubahan. & T-1 \\
\hline
KF-02 & \textit{Feature Serving} Ganda & Sistem mampu melakukan \textit{serving} fitur dalam dua mode, yaitu \textit{online serving} berbasis Valkey melalui protokol RESP untuk \textit{real-time inference} dengan target latensi p99 sub-detik, dan \textit{offline serving} berbasis ClickHouse untuk \textit{historical retrieval} berkapasitas besar yang mendukung \textit{point-in-time correct join} pada data historis. & T-4 \\
\hline
KF-03 & Jaminan Validitas Temporal & Sistem secara bawaan menjamin \textit{point-in-time correctness} ketika menyusun \textit{training dataset} untuk data \textit{time-series}, sehingga mencegah kebocoran data temporal yang dapat membuat evaluasi dan model menjadi tidak valid. & T-3 \\
\hline
KF-04 & Dukungan \textit{Vector Embeddings} & Sistem mendukung penyimpanan, pengindeksan, dan \textit{serving vector embeddings} pada indeks vektor terpisah berbasis Qdrant dengan \textit{similarity search} HNSW melalui antarmuka gRPC dan filter \textit{payload}, dengan target latensi p99 sub-detik untuk volume permintaan menengah. & T-4 \\
\hline
KF-05 & Otomatisasi \textit{Training Pipeline} & Sistem menyediakan \textit{pipeline} otomatis untuk mengorkestrasi proses \textit{model training} dari ujung ke ujung melalui Kubeflow Pipelines, terintegrasi dengan Feast sebagai \textit{feature store}, MLflow untuk \textit{experiment tracking} dan \textit{model registry}, serta Katib untuk \textit{hyperparameter tuning}, sehingga mendukung otomasi setara MLOps \textit{level} 2. & T-3 \\
\hline
KF-06 & Otomatisasi \textit{Deployment} Model & Sistem mengotomatisasi \textit{deployment} model yang telah lulus validasi ke \textit{inference endpoint} produksi melalui KServe \textit{InferenceService} di atas Knative Serving, serta mendukung strategi \textit{rolling update}, \textit{canary deployment}, dan \textit{progressive delivery} berbasis Argo Rollouts dengan mekanisme \textit{rollback} otomatis ketika terjadi degradasi. & T-1 \\
\hline
KF-07 & Pemantauan \textit{Drift} Otomatis & Sistem secara otomatis dan berkala memantau \textit{data drift} dan \textit{concept drift} dengan membandingkan distribusi fitur di produksi terhadap distribusi pada saat pelatihan menggunakan uji statistik (misalnya PSI dan Kolmogorov-Smirnov) dengan ambang adaptif, serta memicu \textit{retraining} otomatis melalui Argo CronWorkflow ketika ambang terlampaui. & T-3 \\
\hline
KF-08 & Tata Kelola \& Penemuan Otomatis & Sistem menangkap dan menyajikan metadata serta \textit{lineage} untuk fitur, \textit{dataset}, model, dan \textit{pipeline} pada katalog DataHub melalui \textit{ingestion} otomatis dan peristiwa \textit{lineage} berbasis OpenLineage, sehingga memudahkan penemuan, analisis dampak, dan pemenuhan kebutuhan \textit{compliance}. & T-2 \\
\hline
KF-09 & Pelacakan Eksperimen Terintegrasi & Sistem menyediakan \textit{experiment tracking} yang terintegrasi untuk mencatat parameter, metrik, artefak, dan versi kode, sehingga proses eksperimen dapat direproduksi dan dibandingkan secara sistematis. & T-2 \\
\hline
KF-10 & \textit{Versioning} dan Registri Model & Sistem menyediakan \textit{model registry} terpusat untuk mengelola versi model dan status siklus hidupnya (\textit{staging}, \textit{production}, \textit{archived}) disertai \textit{audit trail} lengkap untuk mendukung \textit{governance}. & T-2 \\
\hline
KF-11 & Pemrosesan \textit{Batch} dan \textit{Stream} & Sistem mendukung \textit{batch processing} untuk analisis historis dan \textit{feature engineering} serta \textit{stream processing} untuk komputasi fitur \textit{real-time}, dan keduanya terhubung dengan \textit{feature store} sehingga definisi fitur tetap konsisten. & T-4 \\
\hline
KF-12 & Konfigurasi Deklaratif & Sistem memungkinkan konfigurasi komponen seperti fitur, \textit{pipeline}, dan \textit{deployment} dalam format deklaratif (manifes Kubernetes, \textit{overlay} Kustomize, dan definisi Feast berbasis Python) yang dikontrol versinya di Gitea dan disinkronkan oleh Argo CD sebagai \textit{single source of truth}. & T-1 \\
\hline
KF-13 & API dan SDK Terpadu & Sistem menyediakan SDK Python yang seragam (Feast, MLflow, dan klien KServe v2 \textit{inference}) dalam satu \textit{environment} uv sebagai antarmuka klien standar untuk membaca \textit{online feature}, menelusuri \textit{run} dan registri model, serta memanggil \textit{endpoint inference}, sehingga kompleksitas integrasi menurun dan pola pemakaian menjadi konsisten. & T-1 \\
\hline
KF-14 & Manajemen Sumber Daya & Sistem menyediakan kemampuan pengelolaan sumber daya komputasi yang efisien melalui Kueue untuk \textit{job queueing}, prioritas, dan kuota berbasis ClusterQueue, dipadukan dengan HPA, VPA, dan KEDA untuk \textit{dynamic provisioning} pada lingkungan \textit{multi-tenant} agar tidak ada satu \textit{workload} yang menghabiskan sumber daya klaster. & T-1 \\
\hline
\end{longtable}
}